\subsection{Stochastic Variational Baum-Welch}
\label{sec:svi-bw}

For large datasets it is impractical to run the E-step over all training pairs every iteration.
A natural stochastic extension is to sample a minibatch of pairs, run the E-step on those, and update parameters.
Naively replacing the full sufficient statistics with each batch's statistics discards information from previous batches
and results in noisy, slowly converging parameter estimates.

The framework of stochastic variational inference~\cite{hoffman2013svi}
provides a principled resolution in the conditionally conjugate case.
Each M-step parameter group in TKF91\ifdefined\MIXDOMPAPER{} (and TKF92 and MixDom)\else{} (and TKF92, including the mixture-of-TKF92 extension of Section~\ref{sec:em-around-maraschino})\fi{} has an
exponential-family complete-data log-likelihood, regular for the
multinomial groups (mixture weights, fragment-type transitions,
site-class distributions) and curved for the BDI rates and the
reversible-CTMC (GTR) submodel
(Sections~\ref{sec:bdi},~\ref{sec:bridge-expectations}).
For the multinomial groups we use Dirichlet priors that are strictly
conjugate; for the BDI and GTR groups we use independent
Gamma${}\times{}$Dirichlet pseudocounts that are non-conjugate
regularisers in the reversible / curved-family parameterisation
\ifdefined\MIXDOMPAPER (Section~\ref{sec:bw-mixdom-priors})\fi.
For the conjugate groups, Hoffman et al.'s natural-gradient argument
applies in full and the update~\eqref{eq:svi-update} below is exact
natural-gradient ascent on the variational ELBO.
For the non-conjugate groups, the natural-gradient interpretation no
longer holds verbatim, but the same exponential-moving-average update
remains a sound stochastic-MAP-EM scheme: each M-step is affine in the
augmented sufficient statistics, so an EMA of statistics commutes with
the M-step and converges to the MAP fixed point of full-batch EM under
the standard Robbins--Monro step-size conditions.  In practice both
groups are updated by the same rule:
\begin{equation}\label{eq:svi-update}
\bar{T}^{(t+1)} = (1 - \rho_t)\,\bar{T}^{(t)}
  + \rho_t\!\left(\bar{T}_0 + \frac{N}{|B_t|}\,T_{B_t}\right),
\end{equation}
where $\bar{T}^{(t)}$ denotes the running sufficient statistics
at iteration~$t$
(transition counts $\hat{n}_{ij}$,
bridge-expectation dwell times~$W_i$ and jump counts~$U_{ij}$,
character counts~$V_i$,
and BDI expectations $\hat{B}, \hat{D}, \hat{S}$),
$\bar{T}_0$ the prior pseudocounts,
$T_{B_t}$ the batch sufficient statistics from the E-step on minibatch~$B_t$,
$N$ the total number of training pairs,
$|B_t|$ the batch size,
and $\rho_t = (t + \tau)^{-\kappa}$ the Robbins--Monro step size
with delay $\tau \geq 0$ and forgetting rate $\kappa \in (0.5, 1]$.

The M-step---quadratic solve for indel rates~(\ref{eq:kappa-quadratic}),
Dirichlet MAP for mixture weights,
Beta MAP for extension rates,
and the iterative coordinate ascent for $(\exch, \eqm)$---is
then applied to $\bar{T}^{(t+1)}$ exactly as in full-batch EM.
Each M-step maximizes a (penalised) Q-function that is linear in the
sufficient statistics, so averaging the statistics before solving is
equivalent to finding the MAP of the averaged Q-function.
For the conjugate parameter groups this maximizer is also the natural
parameter of the updated variational posterior; for the non-conjugate
groups it is the MAP of the regularised Q-function and not a
variational posterior, but the iteration is identical.

This retains the closed-form exactness of every M-step
while enabling online processing of arbitrarily large pair collections,
with each training pair visited only once.
The two hyperparameters $\tau$ and $\kappa$ have robust defaults
($\tau = 10$, $\kappa = 0.7$) that work across a range of dataset sizes
without tuning~\cite{hoffman2013svi}.

\paragraph{Equivalent formulations of the prior pseudocounts.}
Equation~\eqref{eq:svi-update} places the prior pseudocounts $\bar T_0$
inside the EMA update. An equivalent implementation accumulates only
the data contribution in $\bar T$ and adds $\bar T_0$ at M-step time:
\[
\tilde T^{(t+1)} = (1-\rho_t)\,\tilde T^{(t)} + \rho_t\, \tfrac{N}{|B_t|}\, T_{B_t},
\qquad \theta^{(t+1)} = \mathrm{M\text{-}step}\bigl(\tilde T^{(t+1)} + \bar T_0\bigr).
\]
For every M-step considered in this paper (Dirichlet MAP,
Gamma-augmented counts, pooled GTR formula), the M-step operator is
affine in its sufficient statistics, so
$\mathrm{M}(\bar T^{(t)}) = \mathrm{M}(\tilde T^{(t)} + \bar T_0)$
exactly at every iteration; the two prescriptions yield numerically
identical estimates.
This affineness is what we actually need for the EMA-of-pseudocounts
formulation, and it holds whether or not the prior is strictly
conjugate to the complete-data likelihood
\ifdefined\MIXDOMPAPER(Section~\ref{sec:bw-mixdom-priors} discusses the distinction)\fi. The post-hoc form is slightly more convenient for
implementation because each M-step keeps the prior encapsulated in its
closed-form update; see the code comment at the SVI loop for details.

